\section{Per-City Dataset Statistics and Shadow Taxonomy}
\label{app:data_stats}

This appendix documents the six-category shadow taxonomy referenced in
Section~\ref{sec:annotation} (\S\ref{app:taxonomy}), the per-city and
per-resolution distributional statistics that motivate the LOCO design
(\S\ref{app:perCity}), and the spatial sampling protocol used to draw
patches within each city (\S\ref{app:sampling}).

\subsection{Shadow Taxonomy}
\label{app:taxonomy}

Each annotated shadow segment carries one of six type labels. The taxonomy
is released with the dataset for downstream analyses; the LOCO benchmark
(\S\ref{sec:loco}) collapses these to a binary mask. Definitions and
inclusion rules are in Table~\ref{tab:taxonomy}. Annotators applied a
single inclusion rule across all six types: a region is shadow if it is
solar-directional and consistent with local scene geometry. Sensor
artifacts, dark water, and cloud shadows lacking geometric consistency are
excluded.

\begin{table}[h]
  \centering
  \small
  \setlength{\tabcolsep}{4pt}
  \caption{
    \textbf{Six-category shadow taxonomy.}
    Every annotated segment receives exactly one label.
    The binary LOCO benchmark merges all six into the foreground class.
  }
  \label{tab:taxonomy}
  \begin{tabular}{@{}p{3.0cm}p{9.5cm}@{}}
    \toprule
    \textbf{Category} & \textbf{Definition and inclusion rule} \\
    \midrule
    Building / canyon &
    Cast shadows from buildings and built structures, including narrow
    inter-building canyons. \\
    Under-structure &
    Persistent occlusion shadows beneath elevated structures (bridges,
    overpasses, awnings, carports). \\
    Tree-canopy dapple &
    Shadow patterns cast by tree foliage, including dappled shadow on the
    ground beneath the canopy. \\
    Topography-cast &
    Shadows produced by terrain elevation changes (hills, embankments,
    cuttings). \\
    Vehicle-cast &
    Shadows cast by cars, trucks, and other ground vehicles. \\
    Thin linear &
    Shadows from thin vertical structures (poles, fences, antennas, wires)
    whose width is on the order of a few pixels. \\
    \bottomrule
  \end{tabular}
\end{table}

\subsection{Per-City and Per-Resolution Statistics}
\label{app:perCity}

\begin{figure}[h!]
  \centering
  \includegraphics[width=\linewidth]{figure3.001.jpeg}
  \caption{
    \textbf{\shadow dataset statistics.}
    (A) Sample patches across city-resolution pairs with shadow annotation overlays.
    (B) Shadow-coverage distribution per city (high-res); Phoenix patches are
    predominantly high-coverage ($\ge$30\% shadow area), driving the surface-albedo
    shortcut.
    (C) Median shadow blob size per category and city, showing strong city-level
    variation in shadow morphology.
    (D) Shadow-category distribution per city.
    (E) Mean grayscale intensity distributions per city ($\mu \approx 128$-$132$,
    $\sigma = 15$-$33$); Phoenix's lower standard deviation reflects more uniform
    high-albedo surfaces.
  }
  \label{fig:stats}
\end{figure}

Figure~\ref{fig:stats} reports the distributional properties of the
benchmark. Four observations bear on the LOCO design and on results
reported in Section~\ref{sec:evaluation}.

\textbf{Shadow coverage is city-dependent (panel B).}
Phoenix is dominated by low-coverage patches: $49\%$ of high-resolution
images contain $5$--$15\%$ shadow area and only $9\%$ exceed $30\%$. Miami
and Chicago are shifted toward higher coverage ($50\%$ and $31\%$ of
images above $30\%$ coverage, respectively). Per-image annotation density
follows the same trend ($29.7$ blobs/image in Phoenix, $48.4$ in Miami,
$38.6$ in Chicago).

\textbf{Shadow morphology varies by city (panels C, D).}
Chicago has the largest median tree-canopy blobs ($255$\,px$^2$) and
under-structure blobs ($193$\,px$^2$), reflecting a denser urban grid with
mature street trees. Phoenix shows the smallest under-structure blobs
($72$\,px$^2$) and a vehicle-shadow share of $7\%$, the highest among the
three cities. Miami's distribution is dominated by tree-canopy shadow
($47\%$ of segments) consistent with its dense subtropical canopy. Chicago
is dominated by building shadow ($53\%$).

\textbf{Surface intensity is city-dependent (panel E).}
Mean grayscale intensity distributions overlap in mean
($\mu\!\in\![128, 132]$) but differ in spread: Phoenix
($\sigma\!=\!15$) is roughly half as wide as Miami ($\sigma\!=\!33$). The
narrow Phoenix distribution is consistent with uniform high-albedo desert
surfaces; the wide Miami distribution reflects a mix of bright sand,
asphalt, and dark vegetation.

\textbf{Native resolution affects blob count, not city ranking (panel A).}
The same $384{\times}384$ ground footprint contains substantially more
shadow blobs at $0.6$\,m/px than at $0.3$\,m/px ($+85$ to $+100$ blobs in
the illustrated examples) because the lower-resolution patch covers
$4{\times}$ more area on the ground. The relative ordering across cities is
preserved across resolutions.

These statistics establish that the three cities differ along orthogonal
axes (coverage, morphology, surface intensity) rather than along a single
appearance axis, supporting the LOCO claim in Section~\ref{sec:design}
that the protocol measures transfer across genuine appearance shifts.

\subsection{Spatial Sampling Protocol}
\label{app:sampling}

Patches are drawn from NAIP orthoimagery in two streams: \emph{random}
patches that provide diverse independent samples, and \emph{paired}
patches that share ground coordinates across the two native resolutions.

\textbf{Random patches.}
Within each (city, resolution) pair, target patch counts are distributed
across NAIP tiles overlapping the city administrative boundary. If the
target count is at most the tile count, tiles are sampled without
replacement and each sampled tile contributes one patch; otherwise the
count is distributed evenly with the remainder allocated to randomly
chosen tiles. Within each selected tile, a candidate center is sampled
uniformly inside the tile geometry and accepted only if (i) the point
lies within the city boundary and (ii) a $384{\times}384$ window centered
at the point fits entirely inside the tile (a $192$-pixel buffer to all
edges). Up to $50$ candidates are tried per slot; unfilled slots are
skipped. The one-patch-per-tile rule limits spatial autocorrelation
between samples.

\textbf{Paired patches.}
For each city, $150$ ground coordinates are drawn under the random-patch
procedure on the mid-resolution stream ($0.6$\,m/px), with the additional
constraint that each new center is at least $384$ pixels (one patch
width) from every previously accepted center in the same tile. The
accepted longitude/latitude coordinates are written to a shared registry
and reused for the high-resolution stream ($0.3$\,m/px): for each saved
coordinate, the high-resolution NAIP tile containing the location is
identified and a $384{\times}384$ window is extracted centered on the
exact same geographic point. Each paired patch carries a shared
\texttt{pair\_id} linking its two resolutions. Paired patches support
resolution-matched analyses; results reported in the main paper use the
combined random$+$paired pool unless stated otherwise.

\textbf{Reproducibility.}
Sampling is organized into sequential sessions ($40$ random and $10$
paired patches per resolution per session), each with a distinct random
seed. Per-patch metadata (\texttt{city}, \texttt{scale\_gsd\_m},
\texttt{lon}, \texttt{lat}, \texttt{source\_tile\_id}, \texttt{pair\_id})
is released with the benchmark, enabling exact patch-level reproduction
of the train/val/test splits.